\documentclass[11pt]{article}
\usepackage{mathblatt}
\begin{document}
% Kopfzeile Seite 1 ohne Kennung; die Rückseite (Lösungen, für den Lehrer) bekommt sie nach \begleitteil
\blattkopf{Mathematik}{Basisaufgaben}
\einheitenkopf[][Zettel 17]{Basisaufgaben · Zettel 17}
\zweigzeile{je Aufgabe 1 Punkt · Lösungen auf der Rückseite}
\par\vspace{2pt}\noindent Name \underline{\hspace{7cm}}\hfill Datum \underline{\hspace{3cm}}\hfill Punkte \underline{\hspace{1.2cm}} / 17\par\vspace{6pt}
% \small: volle Seite, kurze Aufgaben zweispaltig (v0.6)
\small

\noindent\begin{minipage}[t]{0.485\linewidth}\raggedright
% 1: Zehnerpotenzschreibweise umwandeln – potenzen-wurzeln-basis-k3-v9 (Original 2025-OS-B1d, ausgedacht: keine Marke)
\begin{aufgabe}{$1\,200 = 1{,}2 \cdot 10^{\square}$ – welche Hochzahl? \leerfeld}
\end{aufgabe}
\end{minipage}\hfill\begin{minipage}[t]{0.485\linewidth}\raggedright
% 2: Größen vergleichen – einheiten-basis-k2-v4 (Original 2026-FOR-B1d, ausgedacht: keine Marke)
\begin{aufgabe}{Vergleiche $45$ mm und $4{,}5$ dm. Setze das richtige Zeichen ein: $<$, $=$ oder $>$. $45$ mm \leerfeld $4{,}5$ dm}
\end{aufgabe}
\end{minipage}\par

\noindent\begin{minipage}[t]{0.485\linewidth}\raggedright
% 3: Zeiteinheiten umrechnen – einheiten-basis-k5-v4 (Original 2025-OS-B1b, ausgedacht: keine Marke)
\begin{aufgabe}{Ein Kuchen backt $1{,}2$ h. Gib die Backzeit in Minuten an. \leerfeld[min]}
\end{aufgabe}
\end{minipage}\hfill\begin{minipage}[t]{0.485\linewidth}\raggedright
% 4: Proportionale Zuordnung Dreisatz – zuordnungen-basis-k2-v4 (Original 2023-OS-B1a, ausgedacht: keine Marke)
\begin{aufgabe}{$3$ m Stoff kosten $21$ €. Wie viel kosten $5$ m? \leerfeld[€]}
\end{aufgabe}
\end{minipage}\par

\noindent\begin{minipage}[t]{0.485\linewidth}\raggedright
% 5: Lineare Gleichung lösen – lineare-gleichungen-basis-k1-v11 (Original 2024-OS-B1d, ausgedacht: keine Marke)
\begin{aufgabe}{Löse die Gleichung $4 \cdot (x - 2{,}5) = 0$. x = \leerfeld}
\end{aufgabe}
\end{minipage}\hfill\begin{minipage}[t]{0.485\linewidth}\raggedright
% 6: Lösung durch Einsetzen prüfen – quadratische-gleichungen-basis-k1-v8 (Original 2025-OS-B1h, ausgedacht: keine Marke)
\begin{aufgabe}{Welcher Wert erfüllt $x \cdot (x + 3) = 10$? \\ \kreuz{$x = -2$} \kreuz{$x = 2$} \kreuz{$x = 5$} \kreuz{$x = 10$}}
\end{aufgabe}
\end{minipage}\par

\noindent\begin{minipage}[t]{0.485\linewidth}\raggedright
% 7: Quadratseite aus Fläche berechnen – flaechen-basis-k2-v3 (Original 2020-OS-B1d, ausgedacht: keine Marke)
\begin{aufgabe}{Ein quadratisches Kissen hat den Flächeninhalt $1600$ cm². Wie lang ist eine Seite? \leerfeld[cm]}
\end{aufgabe}
\end{minipage}\hfill\begin{minipage}[t]{0.485\linewidth}\raggedright
% 8: Median bestimmen – daten-basis-k3-v3 (Original 2024-OS-B1h, ausgedacht: keine Marke)
\begin{aufgabe}{Punkte in einem Test: $18$; $25$; $21$; $16$; $23$; $19$; $22$. Bestimme den Median. \leerfeld Punkte}
\end{aufgabe}
\end{minipage}\par

% 9: Gleichschenkliges Dreieck erkennen – winkel-dreiecke-basis-k2-v3 (Original 2023-OS-B1g, ausgedacht: keine Marke)
\begin{aufgabe}{Im Dreieck ABC ist $\alpha = \beta = 55^\circ$ und AC = $4$ cm. Gib die Länge von BC und die Größe von $\gamma$ an. BC = \leerfeld[cm] , $\gamma =$ \leerfeld °}
\end{aufgabe}

% 10: Grundwert berechnen – prozentrechnung-basis-k1-v5 (Original 2025-OS-B1a, ausgedacht: keine Marke)
\begin{aufgabe}{Beim Kauf eines Fahrradhelms spart Mia $30\,€$, das sind $60\,\%$ des alten Preises. Wie hoch war der alte Preis? \leerfeld[€]}
\end{aufgabe}

% 11: Prozent und Anteil umwandeln – prozentrechnung-basis-k2-v5 (Original 2022-OS-B1f, ausgedacht: keine Marke)
\begin{aufgabe}{Nur $7\,\%$ der Samen keimen nicht. Welche Aussage passt?\\ \kreuz{7 von 10 Samen keimen nicht.}\\ \kreuz{Jeder 7. Samen keimt nicht.}\\ \kreuz{7 von 100 Samen keimen nicht.}}
\end{aufgabe}

% 12: Term zu Sachtext angeben – terme-basis-k2-v10 (Original 2023-OS-B1h, ausgedacht: keine Marke)
\begin{aufgabe}{Das Dreifache einer Zahl $x$ vermindert um 11 ist gleich 25. Welche Gleichung passt? \\ \kreuz{$3 \cdot (x - 11) = 25$} \kreuz{$3x - 11 = 25$} \kreuz{$11 - 3x = 25$} \kreuz{$3x = -11 + 25$}}
\end{aufgabe}

% 13: Figur nach Spiegelung benennen – symmetrie-abbildungen-basis-k1-v4 (Original 2018-OS-B1h, ausgedacht: keine Marke)
\begin{aufgabe}{Ein Rechteck liegt mit einer Seite auf der Geraden $g$ und wird an $g$ gespiegelt. Welche Figur bilden Rechteck und Spiegelbild?}
\end{aufgabe}

% 14: Wahrscheinlichkeit einstufig – bruchrechnung-basis-k1-v7 (Original 2016-OS-B1f, ausgedacht: keine Marke)
\begin{aufgabe}{In einer Lostrommel sind 20 Lose, davon 4 Gewinne. Wie groß ist die Wahrscheinlichkeit, mit einem Los zu gewinnen? \leerfeld}
\end{aufgabe}

% 15: Fehlenden Wert aus Mittelwert bestimmen – daten-basis-k2-v2 (Original 2022-OS-B1d, ausgedacht: keine Marke)
\begin{aufgabe}{Ein Eiscafé verkauft von Montag bis Freitag so viele Becher: $30$; $45$; $40$; \underline{\hspace{1cm}}; $35$. Der Durchschnitt beträgt $38$. Ergänze den fehlenden Wert. \leerfeld Becher}
\end{aufgabe}

% 16: Winkelfunktion Seitenverhältnis angeben – trigonometrie-basis-k2-v7 (Original 2025-OS-B1g, ausgedacht: keine Marke)
\begin{aufgabe}{\begin{minipage}[t]{0.6\linewidth}Rechter Winkel bei $B$. Trage den Bruch ein. $\mathrm{tan}\,\alpha =$ \leerfeld\end{minipage}\hfill\begin{minipage}[t]{0.37\linewidth}\vspace{0pt}
\dreieck{(0,0)}{(3,0)}{(3,2.2)}{r}{s}{t}{\alpha}{}{}
\end{minipage}}
\end{aufgabe}

% 17: Pythagoras Gleichung zuordnen – pythagoras-basis-k1-v2 (Original 2026-FOR-B1j, ausgedacht: keine Marke)
\begin{aufgabe}{\begin{minipage}[t]{0.6\linewidth}Mit welcher Gleichung berechnet man $p$?\\ \kreuz{$p = \sqrt{r^2 + q^2}$}\\ \kreuz{$p = r^2 - q^2$}\\ \kreuz{$p = \sqrt{q^2 - r^2}$}\\ \kreuz{$p = \sqrt{r^2 - q^2}$}\end{minipage}\hfill\begin{minipage}[t]{0.37\linewidth}\vspace{0pt}
\dreieckrw{3.5}{2.5}{p}{q}{r}
\end{minipage}}
\end{aufgabe}

\begleitteil
\blattkopf{Basisaufgaben · Zettel 17}{BAS-Z17 · Lösungen}
\erg{1}{$3$ – das Komma wandert drei Stellen nach links \hfill {\footnotesize\textit{Zehnerpotenzschreibweise umwandeln · 2025}}}
\erg{2}{$45$ mm $= 0{,}45$ dm, also $45$ mm $<$ $4{,}5$ dm \hfill {\footnotesize\textit{Größen vergleichen · 2026}}}
\erg{3}{$1{,}2 \cdot 60 = 72$ min \hfill {\footnotesize\textit{Zeiteinheiten umrechnen · 2025}}}
\erg{4}{Für $1$: $21 : 3 = 7$; $7 \cdot 5 = 35{,}00\,€$ \hfill {\footnotesize\textit{Proportionale Zuordnung Dreisatz · 2023}}}
\erg{5}{$x - 2{,}5 = 0$, also $x = 2{,}5$ \hfill {\footnotesize\textit{Lineare Gleichung lösen · 2024}}}
\erg{6}{$x = 2$, denn $2 \cdot (2 + 3) = 2 \cdot 5 = 10$ \hfill {\footnotesize\textit{Lösung durch Einsetzen prüfen · 2025}}}
\erg{7}{$40 \cdot 40 = 1600$, also $40$ cm \hfill {\footnotesize\textit{Quadratseite aus Fläche berechnen · 2020}}}
\erg{8}{geordnet $16$; $18$; $19$; $21$; $22$; $23$; $25$; der mittlere Wert ist $21$ Punkte \hfill {\footnotesize\textit{Median bestimmen · 2024}}}
\erg{9}{BC = AC = $4$ cm (gleiche Winkel bei A und B); $\gamma = 180^\circ - 2 \cdot 55^\circ = 70^\circ$ \hfill {\footnotesize\textit{Gleichschenkliges Dreieck erkennen · 2023}}}
\erg{10}{$30 : 60 \cdot 100 = 50\,€$ (nicht $60\,\%$ von $30\,€$) \hfill {\footnotesize\textit{Grundwert berechnen · 2025}}}
\erg{11}{7 von 100 Samen keimen nicht. \hfill {\footnotesize\textit{Prozent und Anteil umwandeln · 2022}}}
\erg{12}{$3x - 11 = 25$ \hfill {\footnotesize\textit{Term zu Sachtext angeben · 2023}}}
\erg{13}{Rechteck (doppelt so lang) \hfill {\footnotesize\textit{Figur nach Spiegelung benennen · 2018}}}
\erg{14}{$\frac{4}{20} = \frac{1}{5}$ \hfill {\footnotesize\textit{Wahrscheinlichkeit einstufig · 2016}}}
\erg{15}{$5 \cdot 38 = 190$; $190 - 150 = 40$ \hfill {\footnotesize\textit{Fehlenden Wert aus Mittelwert bestimmen · 2022}}}
\erg{16}{$\mathrm{tan}\,\alpha = \frac{r}{t}$ \hfill {\footnotesize\textit{Winkelfunktion Seitenverhältnis angeben · 2025}}}
\erg{17}{$p = \sqrt{r^2 - q^2}$ \hfill {\footnotesize\textit{Pythagoras Gleichung zuordnen · 2026}}}
\end{document}
